\chapter{Agregación de evidencias por \emph{claim}}\label{cap:agregacion}

\section{Planteamiento}

El clasificador NLI estima
$\hat{p}_k(c, e_i) \approx \Prb(Y = k \mid c, e_i)$ para cada par
\emph{claim}--evidencia. En la práctica, sin embargo, un \emph{claim} $c$
suele evaluarse contra varias evidencias
$\mathcal{E}(c) = \{e_1, \dots, e_n\}$ recuperadas independientemente.
Necesitamos por tanto una regla
$g : (\Delta^{K-1})^n \to \{\textsc{sup}, \textsc{ref}, \textsc{nei},
\textsc{conflict}\}$ que combine las distribuciones por evidencia en un
veredicto único, ampliando el conjunto original con la categoría
\textsc{conflicting} (evidencias que se contradicen entre sí).

\section{Hipótesis simplificadora}

Asumimos \emph{independencia condicional} de las evidencias dado el
veredicto verdadero:
\[
  \Prb(e_1, \dots, e_n \mid Y, c)
  = \prod_{i=1}^n \Prb(e_i \mid Y, c).
\]
Bajo este supuesto y \emph{prior} uniforme $\Prb(Y \mid c) = 1/3$, el
\emph{posterior} agregado se obtiene proporcionalmente a
\[
  \Prb(Y = k \mid c, \mathcal{E}(c))
  \propto \prod_{i=1}^n \hat{p}_k(c, e_i).
\]
La hipótesis de independencia condicional es claramente una
aproximación: en la práctica las evidencias recuperadas pueden
solaparse o estar correlacionadas. La discusión de las desviaciones se
incluye en el anexo.

\section{Regla de decisión propuesta}

Sea $\bar{p}_k$ la probabilidad agregada normalizada y sea
$\tau_{\text{conf}} \in (0, 1)$ un umbral de confianza (en este trabajo
$\tau_{\text{conf}} = 0.55$, calibrado en \emph{dev}). Definimos:

\begin{enumerate}[leftmargin=*,itemsep=0.3em]
  \item Sea $S = \{i : \arg\max_k \hat{p}_k(c, e_i) = \textsc{sup}\}$,
        $R = \{i : \arg\max_k \hat{p}_k(c, e_i) = \textsc{ref}\}$.
  \item \textbf{Conflicto:} si $|S| \geq 1$ y $|R| \geq 1$ y la
        confianza media en ambos grupos supera $\tau_{\text{conf}}$,
        \[
          g(\dots) = \textsc{conflicting}.
        \]
  \item \textbf{Soportada / Refutada:} en caso contrario, si
        $\max_k \bar{p}_k \geq \tau_{\text{conf}}$,
        \[
          g(\dots) = \arg\max_k \bar{p}_k \in
          \{\textsc{sup}, \textsc{ref}\}.
        \]
  \item \textbf{NEI:} si ninguna de las condiciones anteriores se
        cumple,
        \[
          g(\dots) = \textsc{nei}.
        \]
\end{enumerate}

\section{Propiedades}

\paragraph{Consistencia.} Si todas las evidencias coinciden con
confianza alta, $g$ devuelve la etiqueta común. Si el conjunto está
vacío, $g$ devuelve \textsc{nei} por convención.

\paragraph{Manejo de incertidumbre.} El umbral $\tau_{\text{conf}}$
controla el \emph{trade-off} entre cobertura
(proporción de \emph{claims} con veredicto $\neq$ \textsc{nei}) y
precisión. La curva precisión--cobertura para distintos
$\tau_{\text{conf}}$ se reporta en el capítulo~\ref{cap:resultados}.

\paragraph{Detección de contradicciones.} La introducción de la
categoría \textsc{conflicting} es una contribución funcional del
sistema: en la formulación original FEVER las evidencias contradictorias
se resuelven por mayoría y se pierde la señal de inconsistencia.

\section{Implementación}

La regla está implementada en
\path{fakenews-backend/fever/aggregation.py} (función
\texttt{aggregate\_claim\_predictions}) y validada por las pruebas
unitarias en \path{fakenews-backend/tests/}. La implementación es
puramente funcional y no requiere parámetros entrenables más allá del
umbral $\tau_{\text{conf}}$.
